%=====================================================================
%  Criptografía y Seguridad - ITBA
%  Clase 1 — Introducción a la Criptografía
%  Apunte integral: teoría (esta clase NO tuvo clase práctica asociada)
%=====================================================================
\documentclass[11pt,a4paper]{article}

%---------------------------------------------------------------------
%  Paquetes
%---------------------------------------------------------------------
\usepackage{iftex}
\ifPDFTeX
  \usepackage[utf8]{inputenc}
  \usepackage[T1]{fontenc}
  \usepackage{lmodern}
\else
  \usepackage{fontspec}
\fi
\usepackage[spanish,es-tabla,es-noshorthands]{babel}
\usepackage{microtype}

\usepackage{amsmath,amssymb,amsthm}
\usepackage{mathtools}
\usepackage{geometry}
\geometry{a4paper,margin=2.4cm,headheight=22pt}

\usepackage{xcolor}
\usepackage{graphicx}
\usepackage{enumitem}
\usepackage{booktabs}
\usepackage{array}
\usepackage{tcolorbox}
\tcbuselibrary{skins,breakable}
\usepackage{fancyhdr}
\usepackage{titlesec}
\usepackage{hyperref}
\usepackage{listings}
\usepackage{caption}

\usepackage{tikz}
\usetikzlibrary{shapes.geometric,shapes.misc,shapes.symbols,arrows.meta,positioning,calc,
                fit,backgrounds,decorations.pathmorphing,shadows.blur,chains}

%---------------------------------------------------------------------
%  Paleta de color (inspirada en las diapositivas de la cátedra)
%---------------------------------------------------------------------
\definecolor{itbaCyan}{HTML}{6EC6D9}
\definecolor{itbaCyanDark}{HTML}{2E8B9E}
\definecolor{itbaBlue}{HTML}{AEDCEA}
\definecolor{boxBlue}{HTML}{BBD4F0}
\definecolor{slideGray}{HTML}{ECECEC}
\definecolor{slideCream}{HTML}{F4F1DE}
\definecolor{practGreen}{HTML}{4F8A3D}
\definecolor{practGreenL}{HTML}{E2EFD9}
\definecolor{attackRed}{HTML}{C0392B}
\definecolor{codeBg}{HTML}{F5F5F5}
\definecolor{histBar}{HTML}{2E8B9E}
\definecolor{timelineBlue}{HTML}{1F4E79}

%---------------------------------------------------------------------
%  Hipervínculos
%---------------------------------------------------------------------
\hypersetup{
  colorlinks=true,
  linkcolor=itbaCyanDark,
  urlcolor=itbaCyanDark,
  citecolor=itbaCyanDark,
  pdftitle={Clase 1 - Introduccion a la Criptografia},
  pdfauthor={Criptografia y Seguridad - ITBA}
}

%---------------------------------------------------------------------
%  Encabezados y pies
%---------------------------------------------------------------------
\pagestyle{fancy}
\fancyhf{}
\lhead{\small\textsc{Criptografía y Seguridad — ITBA}}
\rhead{\small Clase 1 · Introducción a la Criptografía}
\cfoot{\small\thepage}
\renewcommand{\headrulewidth}{0.4pt}
\renewcommand{\headrule}{\hbox to\headwidth{\color{itbaCyanDark}\leaders\hrule height \headrulewidth\hfill}}

%---------------------------------------------------------------------
%  Títulos de sección con barra de color (estilo diapositiva)
%---------------------------------------------------------------------
\titleformat{\section}
  {\normalfont\Large\bfseries\color{itbaCyanDark}}
  {\thesection}{0.6em}{}
  [\vspace{-0.4em}{\color{itbaCyan}\titlerule[2pt]}]
\titleformat{\subsection}
  {\normalfont\large\bfseries\color{itbaCyanDark}}
  {\thesubsection}{0.6em}{}
\titleformat{\subsubsection}
  {\normalfont\normalsize\bfseries\color{black}}
  {\thesubsubsection}{0.6em}{}

%---------------------------------------------------------------------
%  Cajas reutilizables
%---------------------------------------------------------------------
% Caja "diapositiva" (recuadro gris de las slides)
\newtcolorbox{slidebox}[1][]{
  enhanced, colback=slideGray, colframe=black!55, boxrule=0.5pt,
  arc=1pt, left=6pt,right=6pt,top=4pt,bottom=4pt, #1}

% Caja "aclaración" (cuadro crema con globo de las slides)
\newtcolorbox{notebox}[1][]{
  enhanced, breakable, colback=slideCream, colframe=black!45, boxrule=0.6pt,
  arc=3pt, left=7pt,right=7pt,top=5pt,bottom=5pt, #1}

% Caja de definición
\newtcolorbox{defbox}[1]{
  enhanced, breakable, colback=itbaBlue!25, colframe=itbaCyanDark,
  boxrule=1pt, arc=2pt, left=8pt,right=8pt,top=6pt,bottom=6pt,
  fonttitle=\bfseries, coltitle=white, title={#1}}

% Caja de nota / ejemplo trabajado
\newtcolorbox{practbox}[1]{
  enhanced, breakable, colback=practGreenL, colframe=practGreen,
  boxrule=1pt, arc=2pt, left=8pt,right=8pt,top=6pt,bottom=6pt,
  fonttitle=\bfseries, coltitle=white, title={#1}}

% Caja de atención / ataque
\newtcolorbox{warnbox}[1]{
  enhanced, breakable, colback=attackRed!8, colframe=attackRed,
  boxrule=1pt, arc=2pt, left=8pt,right=8pt,top=6pt,bottom=6pt,
  fonttitle=\bfseries, coltitle=white, title={#1}}

%---------------------------------------------------------------------
%  Estilo de código / monoespaciado
%---------------------------------------------------------------------
\lstdefinestyle{mono}{
  basicstyle=\ttfamily\footnotesize,
  backgroundcolor=\color{codeBg},
  breaklines=true, frame=single, rulecolor=\color{black!30},
  showstringspaces=false, columns=fullflexible, tabsize=2,
  xleftmargin=10pt, framexleftmargin=8pt}

%---------------------------------------------------------------------
%  Macros matemáticos
%---------------------------------------------------------------------
\newcommand{\Pset}{\mathcal{P}}
\newcommand{\Cset}{\mathcal{C}}
\newcommand{\Kset}{\mathcal{K}}
\newcommand{\Mset}{\mathcal{M}}
\DeclareMathOperator{\Enc}{Enc}
\DeclareMathOperator{\Dec}{Dec}
\DeclareMathOperator{\Gen}{Gen}
\DeclareMathOperator{\mcd}{mcd}

%=====================================================================
\begin{document}
%=====================================================================

%---------------------------------------------------------------------
%  Portada
%---------------------------------------------------------------------
\begin{titlepage}
\centering
\vspace*{1.0cm}

\begin{tikzpicture}
  \node[rounded corners=2pt, fill=itbaCyan, text=black, inner sep=14pt,
        minimum width=0.8\textwidth, align=center] (t)
        {\Huge\bfseries Criptografía y Seguridad};
\end{tikzpicture}

\vspace{0.6cm}
{\LARGE Primera mitad: Criptografía}\\[0.25cm]
{\Huge\bfseries\color{itbaCyanDark} Introducción a la Criptografía}

\vspace{0.9cm}

% Ilustración: bóveda / caja fuerte (recreación del icono de portada)
\begin{tikzpicture}[scale=1.0]
  % marco de la bóveda
  \fill[black!12] (-3.2,-2.4) rectangle (3.2,2.8);
  \draw[line width=2pt, black!55] (-3.2,-2.4) rectangle (3.2,2.8);
  % puerta blindada
  \fill[itbaCyanDark!75] (-2.4,-2.0) rectangle (2.4,2.4);
  \draw[line width=3pt, black!70] (-2.4,-2.0) rectangle (2.4,2.4);
  \draw[line width=1.2pt, itbaCyan] (-2.0,-1.6) rectangle (2.0,2.0);
  % bisagra central
  \draw[line width=2pt, black!70] (0,-2.0) -- (0,2.4);
  % volantes / ruedas
  \foreach \a in {0,45,...,315}{
    \draw[line width=2pt, black!70] (-1.1,0.2) -- ++(\a:0.6);
    \draw[line width=2pt, black!70] ( 1.1,0.2) -- ++(\a:0.6);
  }
  \draw[line width=2.5pt, black!70] (-1.1,0.2) circle (0.42);
  \fill[itbaCyan] (-1.1,0.2) circle (0.16);
  \draw[line width=2.5pt, black!70] (1.1,0.2) circle (0.42);
  \fill[itbaCyan] (1.1,0.2) circle (0.16);
  % pernos
  \foreach \y in {-1.2,-0.4,0.8,1.6}{
     \fill[black!70] (-2.2,\y) circle (0.07);
     \fill[black!70] (2.2,\y) circle (0.07);
  }
\end{tikzpicture}

\vfill
{\large Apunte integral — Teoría}\\[0.2cm]
{\normalsize Clase 1 \quad·\quad Capítulo 1, \emph{Introduction to Modern Cryptography} (Katz \& Lindell)}\\[0.4cm]
{\small Instituto Tecnológico de Buenos Aires (ITBA)}

\vspace{0.5cm}
{\footnotesize\itshape Documento elaborado a partir de la transcripción de la clase teórica
(Prof.\ Pablo E.\ Abad) y de las diapositivas oficiales.}

\vspace{0.4cm}
\begin{tcolorbox}[enhanced,colback=practGreenL,colframe=practGreen,boxrule=1pt,
  arc=3pt,width=0.8\textwidth,halign=center]
\small\textbf{Nota.} Esta clase \textbf{no tuvo clase práctica asociada}: fue exclusivamente
teórica. (El ejercicio de descifrado se realizó en vivo durante la propia clase.)
\end{tcolorbox}
\end{titlepage}

\tableofcontents
\newpage

%=====================================================================
\section{¿Qué es la criptografía?}
%=====================================================================

Esta es la primera clase de la mitad de la materia dedicada a \textbf{criptografía}. El
enfoque es introductorio: se trabajan \emph{fundamentos e ideas}, sin entrar todavía en el
campo formal ni en construcciones de uso actual. La estrategia pedagógica es
\textbf{aprovechar la historia de la criptografía} para ir entendiendo los conceptos básicos.

\begin{defbox}{Criptografía}
\begin{itemize}[leftmargin=1.4em,itemsep=2pt]
  \item \textbf{Etimología}: del griego, significa \emph{escritura secreta}. Sus orígenes
        responden a la necesidad de \textbf{ocultar información}.
  \item \textbf{Definición moderna}: un \emph{conjunto de funciones matemáticas y técnicas}
        que permiten asegurar distintas \emph{propiedades} sobre la información.
  \item \textbf{Rol en seguridad}: son las \emph{herramientas básicas} desde donde se
        construye la seguridad.
\end{itemize}
\end{defbox}

El concepto moderno trasciende la mera \textbf{confidencialidad} (esconder información) y
abarca otras propiedades. Pero la idea nace con un hecho atávico: \emph{desde que el hombre
se agrupa, existe la necesidad de mantener algo en secreto}.

\subsection{La criptografía está en todos lados (y oculta)}
Fuera del mundo técnico casi no se habla de criptografía, pero sus construcciones están
prácticamente en todo lo que usamos todo el tiempo. Ejemplos mencionados:
\begin{itemize}[leftmargin=1.4em,itemsep=1pt]
  \item \textbf{Comunicaciones seguras}: tráfico web, tráfico inalámbrico (Wi-Fi),
        transmisiones satelitales, tráfico por Internet (cableado o no).
  \item \textbf{Protección de archivos en disco}: si a un empleado le roban o pierde la
        notebook, los datos confidenciales de la empresa deben quedar protegidos.
  \item \textbf{Autenticación de usuarios}: segundo factor de autenticación, generadores de
        claves, Google Authenticator.
  \item \textbf{Protección de contenido} (DRM): venta de libros digitales (e-readers),
        alquiler de películas en formato digital para evitar la copia indiscriminada.
  \item \textbf{Firmas digitales}: contratos que se celebran de forma totalmente digital, sin
        papel.
  \item \textbf{Voto electrónico}: con más o menos sofisticación, hay criptografía detrás.
  \item \textbf{Dinero electrónico}: desde Bitcoin hasta todo el ecosistema de criptomonedas.
\end{itemize}

%=====================================================================
\section{El criptosistema (visión informal)}
%=====================================================================

Dentro de la larga lista de funciones criptográficas, hay una construcción \textbf{básica,
común y muy ligada al origen}: el \textbf{criptosistema}. Lo abordamos primero de forma
intuitiva; la definición formal llega más adelante.

\begin{defbox}{Criptosistema (idea inicial)}
Es, en principio, un conjunto de \textbf{dos funciones}:
\begin{itemize}[leftmargin=1.4em,itemsep=2pt]
  \item \textbf{Cifrado} (cifrar / encriptar): toma un mensaje \emph{arbitrario} (un texto,
        un archivo, cualquier cosa) que tiene sentido para nosotros y lo transforma en algo
        que parece producido por un mono tipeando teclas al azar.
  \item \textbf{Descifrado} (descifrar / desencriptar): toma ese mensaje sin sentido (el
        \emph{mensaje cifrado}) y recupera el \textbf{mensaje original}.
\end{itemize}
\end{defbox}

% ---- Diagrama: esquema de comunicación cifrada (slide Criptosistema) ----
\begin{figure}[h]
\centering
\begin{tikzpicture}[font=\small,>=Stealth,node distance=8mm]
  \tikzset{
    plano/.style={draw,rounded corners,fill=slideCream,minimum width=2.0cm,
                  minimum height=1.3cm,align=center},
    cif/.style={draw,rounded corners,fill=boxBlue,minimum width=2.0cm,
                minimum height=1.3cm,align=center},
    gear/.style={draw,circle,fill=black!80,text=white,minimum size=0.85cm,
                 font=\scriptsize\bfseries}
  }
  \node[plano] (m1) {Mensaje o\\ texto plano};
  \node[gear, right=1.0cm of m1] (e) {E/D};
  \node[cif, right=1.0cm of e] (c) {Mensaje o\\ texto cifrado};
  \node[gear, right=1.0cm of c] (d) {E/D};
  \node[plano, right=1.0cm of d] (m2) {Mensaje o\\ texto plano};

  \draw[->,line width=1pt] (m1) -- (e);
  \draw[->,line width=1pt] (e) -- (c);
  \draw[->,line width=1pt] (c) -- (d);
  \draw[->,line width=1pt] (d) -- (m2);

  % clave
  \node[draw,fill=yellow!75!orange,rounded corners,minimum width=1.0cm,
        minimum height=0.5cm,font=\scriptsize] (k) at ($(c)+(0,-2.0)$) {clave $k$};
  \draw[->,line width=1pt] (k) -| (e);
  \draw[->,line width=1pt] (k) -| (d);

  % etiquetas globo
  \node[draw,rounded corners,fill=practGreenL,draw=practGreen,font=\scriptsize,
        above=4mm of e] {Cifrar / encriptar};
  \node[draw,rounded corners,fill=practGreenL,draw=practGreen,font=\scriptsize,
        above=4mm of d] {Descifrar / desencriptar};

  % ejemplo
  \node[cif,below=2.7cm of m1,font=\scriptsize] (ex1) {Nos vemos\\ a las 18};
  \node[cif,below=2.7cm of c,font=\scriptsize]  (ex2) {skjejcaeidl\\ dsfjkfj ieu\\ olkjds kjkli};
  \node[cif,below=2.7cm of m2,font=\scriptsize] (ex3) {Nos vemos\\ a las 18};
  \node[left=1mm of ex1,font=\scriptsize\itshape] {Ejemplo:};
\end{tikzpicture}
\caption{Esquema de comunicación cifrada. El emisor cifra el mensaje plano con la clave; el
mensaje cifrado puede viajar por un canal inseguro; el receptor lo descifra con la misma
clave. \textit{(Recreación de la diapositiva ``Criptosistema''.)}}
\end{figure}

\subsection{La función de cifrado y la clave}
La función de cifrado es, conceptualmente, una función de \textbf{dos parámetros}: el mensaje
a transformar y un bloque adicional de información llamado \textbf{clave}.
\[
  \underbrace{\text{mensaje plano}}_{p}\ \times\ \underbrace{\text{clave}}_{k}
  \;\longrightarrow\; \underbrace{\text{mensaje cifrado}}_{c}
  \qquad\Longrightarrow\qquad e_k(p)=c
\]

\begin{notebox}
\textbf{Diferentes representaciones, mismo significado.} En la literatura aparece la misma
función de cifrado escrita de varias maneras:
\[
  e_k(p) \;=\; \mathrm{enc}_k(p) \;=\; e(k,p) \;=\; \{\,p\,\}_k
\]
La última notación, $\{\,p\,\}_k$ (``lo que está entre llaves está cifrado con la clave $k$''),
se usa más en \emph{protocolos}. Análogamente, el descifrado se escribe
$d_k(c)=\mathrm{dec}_k(c)=d(k,c)=\{\,c\,\}_k^{-1}$ y cumple
\[
  \text{mensaje cifrado}\ \times\ \text{clave} \;\longrightarrow\; \text{mensaje plano},
  \qquad d(c,k)=p .
\]
\end{notebox}

\begin{defbox}{¿Qué es la clave?}
La clave es \textbf{un bloque arbitrario de información}, igual que el mensaje. La separamos
y le damos un nombre especial porque, al probar la seguridad de un sistema, vamos a
\textbf{presuponer que la clave es información protegida}: el atacante \emph{no} la tiene, y
si la consigue, se acabó la seguridad.
\end{defbox}

\begin{notebox}
\textbf{Ejercicio mental del atacante.} Una técnica habitual para razonar sobre seguridad es
\emph{ponerse en el lugar del atacante}: ``si yo quisiera romper esto, ¿qué tendría que
pasar?''. En todos esos escenarios la constante es: donde aparece una clave, suponemos que el
atacante no la tiene.
\end{notebox}

\subsection{Principio de Kerckhoffs}
\begin{defbox}{Principio de Kerckhoffs}
Un criptosistema debe ser seguro \textbf{incluso si todo el sistema} —todo el algoritmo,
todas las funciones— \textbf{excepto la clave es de conocimiento público.}
\end{defbox}
\noindent En consecuencia, la seguridad de un sistema \textbf{no puede depender} de que el
atacante desconozca, por ejemplo, el código fuente. Debemos asumir que eventualmente lo
conoce, y eso no debe comprometer la seguridad.

\subsection{¿Para qué sirve un criptosistema?}
En el fondo, resuelve el problema de \textbf{controlar el acceso y la diseminación de
información}. Proteger un documento ``a mano'' (borrarlo de todos lados salvo de un pendrive
guardado en una caja fuerte) no escala: en cuanto el archivo está digitalizado y hay que
compartirlo, se pierde el control sobre quién accede.

\begin{notebox}
\textbf{La gran idea: ``patear'' el problema a algo manejable.} Proteger un mensaje de tamaño
arbitrario es difícil. Con un criptosistema, ciframos el mensaje con una clave: el
\textbf{mensaje cifrado} puede dejarse en cualquier lado (una página pública, enviarlo por
todos lados) siempre que la \textbf{clave esté protegida}. No resolvemos el problema por
completo, pero lo reducimos a proteger algo \textbf{mucho más chico}: en los criptosistemas
modernos, la clave no tiene relación con el tamaño del mensaje. Se pueden cifrar
\emph{terabytes} de información con una clave de \textbf{128 bits}, y proteger 128 bits es
mucho más fácil que proteger varios terabytes.
\end{notebox}

% ---- Diagrama: uso básico (slide Uso básico de un criptosistema) ----
\begin{figure}[h]
\centering
\begin{tikzpicture}[font=\small,>=Stealth]
  \tikzset{
    cif/.style={draw,rounded corners,fill=boxBlue,minimum width=1.9cm,
                minimum height=1.2cm,align=center,font=\scriptsize},
    gear/.style={draw,circle,fill=black!80,text=white,minimum size=0.8cm,
                 font=\scriptsize\bfseries}
  }
  % lado izquierdo: sin clave
  \node[font=\bfseries] at (-3.4,1.3) {Sin conocer la clave:};
  \node[cif] (c1) at (-5.0,0) {skjejcaeidl\\ dsfjkfj ieu};
  \node[gear,right=0.7cm of c1] (g1) {E/D};
  \node[cif,right=0.7cm of g1] (o1) {hsdfudsahj\\ sjhguy};
  \draw[->] (c1)--(g1); \draw[->] (g1)--(o1);
  \node[draw,fill=yellow!75!orange,rounded corners,minimum width=0.9cm,
        minimum height=0.45cm,font=\scriptsize] (k1) at ($(g1)+(0,-1.2)$) {clave};
  \node[font=\Large,red] at (k1) {$\times$};
  \draw[->] (k1) -- (g1);

  % separador
  \draw[line width=1pt,black!50] (0.2,-1.8) -- (0.2,1.5);

  % lado derecho: con clave
  \node[font=\bfseries] at (4.0,1.3) {Conociendo la clave:};
  \node[cif] (c2) at (2.4,0) {skjejcaeidl\\ dsfjkfj ieu};
  \node[gear,right=0.7cm of c2] (g2) {E/D};
  \node[cif,right=0.7cm of g2] (o2) {Nos vemos\\ a las 18};
  \draw[->] (c2)--(g2); \draw[->] (g2)--(o2);
  \node[draw,fill=yellow!75!orange,rounded corners,minimum width=0.9cm,
        minimum height=0.45cm,font=\scriptsize] (k2) at ($(g2)+(0,-1.2)$) {clave};
  \draw[->] (k2) -- (g2);
\end{tikzpicture}
\caption{El texto cifrado \emph{no tiene información útil} y puede ser almacenado o
transferido libremente. Sin la clave no se recupera el mensaje; con la clave, sí. Permite
\textbf{controlar quién puede recuperar un mensaje}. \textit{(Recreación de la diapositiva
``Uso básico de un criptosistema''.)}}
\end{figure}

%=====================================================================
\section{Seguridad informal de un criptosistema}
%=====================================================================

Una primera respuesta intuitiva (aportada en clase) es: \emph{``alguien que no tiene la clave
no puede descifrar el mensaje en un tiempo razonable''}. Es correcta, pero la criptografía
moderna pide \textbf{más}.

\begin{defbox}{Seguridad (informal)}
Un criptosistema es seguro si \textbf{ningún adversario puede computar cualquier función del
texto plano} a partir del mensaje cifrado que posee. Es decir, no se puede extraer
\emph{ningún tipo} de información del mensaje. Significa:
\begin{itemize}[leftmargin=1.4em,itemsep=1pt]
  \item No se puede recuperar el mensaje.
  \item No se puede recuperar \emph{parte} del mensaje.
  \item No se puede recuperar el \emph{sentido} del mensaje.
\end{itemize}
\end{defbox}

\begin{notebox}
\textbf{¿Por qué pedir tanto?} A veces un adversario no necesita todo el documento:
\begin{itemize}[leftmargin=1.3em,itemsep=1pt]
  \item Un informe gigante puede tener un \emph{resumen} en la primera página: recuperar solo
        esa página ya hace daño.
  \item En la negociación secreta de una fusión multimillonaria, a quien espía (para invertir
        ilegalmente en la bolsa) no le hace falta decodificar el mensaje: le basta con saber
        si el resultado será \emph{positivo o negativo}.
\end{itemize}
Por eso la seguridad se asocia a que \textbf{no se pueda extraer ningún tipo de información}.
\end{notebox}

\subsection{Criptografía vs.\ Criptoanálisis}
\begin{notebox}
Las técnicas para probar o romper la seguridad de un criptosistema se denominan
\textbf{criptoanálisis}. La disciplina suele separarse en:
\begin{itemize}[leftmargin=1.3em,itemsep=1pt]
  \item \textbf{Criptografía}: el estudio (diseño) de las funciones.
  \item \textbf{Criptoanálisis}: el estudio de los ataques.
\end{itemize}
En esta materia se verá criptoanálisis ``muy por arriba'', para saber que existe.
\end{notebox}

\begin{warnbox}{No inventes tu propia primitiva}
Un objetivo secundario de la materia es que el ingeniero entienda la \textbf{complejidad}
detrás de las funciones criptográficas, para que \textbf{no} se le ocurra crear una primitiva
nueva y, en cambio, \emph{busque las que ya existen}. En los años 80 diseñar funciones
criptográficas era un pasatiempo de ingenieros; hoy, dado el avance del criptoanálisis, se
estima que hay en el planeta \textbf{solo unas 30 o 40 personas} capacitadas para construir
funciones que sobrevivan al escrutinio público (y cada una con su colección de fracasos). Por
suerte, como ingenieros vamos a \emph{utilizar} las que ya existen.
\end{warnbox}

%=====================================================================
\section{Una breve historia de la criptografía}
%=====================================================================

La necesidad de ocultar información existe desde siempre. La siguiente línea de tiempo guía el
recorrido por los criptosistemas clásicos.

% ---- Diagrama: línea de tiempo (slide Historia de la criptografía) ----
\begin{figure}[h]
\centering
\begin{tikzpicture}[font=\footnotesize]
  \def\xl{0}      % x de la línea
  \def\xt{0.9}    % x del texto
  % barra vertical
  \fill[timelineBlue] (\xl-0.12,-13.2) rectangle (\xl+0.12,0.3);
  \foreach \y/\fecha/\titulo/\desc in {
     0/{$\sim$2500 a.C.}/{Egipto}/{Criptosistemas de sustitución: jeroglíficos no estándar},
     -1.9/{$\sim$700--300 a.C.}/{Escítala (Scitala)}/{Criptosistema de transposición usando báculos},
     -3.8/{50 a.C.}/{César}/{Criptosistema de rotación (subtipo de sustitución)},
     -5.7/{800 d.C.}/{Primeros documentos de criptoanálisis}/{Análisis de frecuencias: comienzan a ``romperse'' los cifrados},
     -7.6/{1500}/{Criptosistemas polialfabéticos}/{Un símbolo cifrado representa distintos símbolos del original},
     -9.5/{1939}/{Fuerza bruta: Bombe -- Enigma}/{Ataques sistemáticos por fuerza bruta -- protocomputadoras},
     -11.4/{1949}/{Information Theory \& Cryptography (Shannon)}/{Publicaciones seminales: nace la criptografía moderna}
  }{
     \fill[timelineBlue] (\xl,\y) circle (0.13);
     \draw[timelineBlue,line width=1pt] (\xl-0.55,\y) -- (\xl,\y);
     \node[anchor=east,timelineBlue] at (\xl-0.6,\y) {\fecha};
     \node[anchor=west,timelineBlue,font=\footnotesize\bfseries] at (\xt,\y+0.18) {\titulo};
     \node[anchor=west,text width=8.6cm,timelineBlue] at (\xt,\y-0.32) {\desc};
  }
\end{tikzpicture}
\caption{Línea de tiempo de la criptografía. \textit{(Recreación de la diapositiva ``Historia
de la criptografía''.)}}
\end{figure}

\begin{itemize}[leftmargin=1.4em,itemsep=2pt]
  \item \textbf{Egipto ($\sim$2500 a.C.)}: indicios de criptosistemas de \emph{sustitución}
        basados en darle un significado distinto a los jeroglíficos. (Hay registros hebreos
        de hasta $\sim$5000 a.C.\ que se presumen ``proto-criptografía''.)
  \item \textbf{Escítala (época grecorromana)}: criptosistema de \emph{transposición}. Se
        enrolla una tira de cuero en un báculo de cierto tamaño y topología, se escribe
        \emph{transversalmente} y luego se desenrolla. Para descifrar hace falta un báculo
        idéntico. Lo usaban los gobernantes de las polis griegas.
  \item \textbf{Julio César}: usos documentados de criptografía en campañas militares para
        coordinar ataques entre unidades. En su honor existe el \emph{cifrado de César}.
  \item \textbf{$\sim$800 d.C.}: primeros intentos documentados de \emph{romper}
        criptosistemas (análisis de frecuencias).
  \item \textbf{Renacimiento ($\sim$1500)}: surge una generación nueva, los
        \emph{polialfabéticos}, que dominan casi \textbf{400 años}. Primera gran revolución.
  \item \textbf{Segunda Guerra Mundial}: segunda gran revolución. Uso bélico masivo de
        criptografía y, por primera vez, un esfuerzo deliberado y muy financiado por
        \emph{quebrarla} en todos los frentes.
  \item \textbf{Shannon (1949)}: funda la \emph{Teoría de la Información} y da origen a la
        \textbf{criptografía moderna}. Ochenta años después se siguen usando sus fundamentos
        teóricos (llevados al poder de cómputo actual).
\end{itemize}

%=====================================================================
\section{Cifrado por rotación (ROT-X) y el cifrado de César}
%=====================================================================

Es el criptosistema más básico que se conoce.

\begin{defbox}{Cifrado por rotación}
\begin{itemize}[leftmargin=1.4em,itemsep=2pt]
  \item Se asigna un orden a los símbolos del alfabeto (cada letra es un número: su posición,
        empezando por 0).
  \item La \textbf{clave} $k$ es el \emph{número de desplazamientos}; varía entre $1$ y $N$
        (cantidad de letras). En realidad entre $1$ y $N-1$, porque rotar el alfabeto completo
        da la función identidad.
  \item \textbf{Cifrar}: reemplazar cada letra por la que está $k$ posiciones después, de
        forma \textbf{cíclica} (después de la \texttt{z} se vuelve a la \texttt{a}).
\end{itemize}
\end{defbox}

\begin{practbox}{Ejemplo trabajado: $e(\text{prueba},4)$}
Con clave $k=4$ se reemplaza cada letra por la que está 4 posiciones después:
\begin{center}
\renewcommand{\arraystretch}{1.2}
\begin{tabular}{cccccc}
\toprule
plano   & p & r & u & e & b a\\
\midrule
\textbf{cifrado} & t & v & y & i & f\,$\cdots$\\
\bottomrule
\end{tabular}
\end{center}
Resultado: $\;e(\text{prueba},4)=\texttt{tvyife}$. El cifrado preserva el largo del mensaje.
\end{practbox}

\subsection{Criptoanálisis: fuerza bruta}
\begin{warnbox}{Debilidad: hay pocas claves}
Como el atacante \emph{conoce el algoritmo} (principio de Kerckhoffs) y \emph{desconoce la
clave}, no hay forma de impedir un \textbf{ataque por fuerza bruta}: probar todas las claves.
Con tan pocas claves posibles, el ataque es trivial.
\end{warnbox}

\begin{practbox}{Ataque por fuerza bruta sobre $c=\texttt{tvyife}$}
Se prueban todos los desplazamientos hasta hallar uno ``con sentido'':
\begin{center}
\renewcommand{\arraystretch}{1.15}
\begin{tabular}{lll}
\toprule
clave & $d(k,\texttt{tvyife})$ & \\
\midrule
$k=1$ & \texttt{suxhed} & $\leftarrow$ sin sentido\\
$k=2$ & \texttt{rtwgdc} & $\leftarrow$ sin sentido\\
$k=3$ & \texttt{qsvfcb} & $\leftarrow$ sin sentido\\
$k=4$ & \texttt{prueba} & $\leftarrow$ \textbf{¡con sentido!}\\
\bottomrule
\end{tabular}
\end{center}
Resultado: $k=4$, \ $m=\texttt{prueba}$.
\end{practbox}

\begin{notebox}
\textbf{El cifrado de César concreto.} El que usó Julio César y está documentado para
coordinar campañas bélicas era una rotación \emph{fija} de \textbf{3 letras} en el alfabeto
romano. Inseguro en retrospectiva, pero le valió un montón de guerras (en su época, las
condiciones eran otras).
\end{notebox}

\subsection{Redundancia, distancia de unicidad y los límites de la fuerza bruta}
Una buena pregunta surge naturalmente: ¿y si al probar con otra clave también aparece una
palabra con sentido? Esta inquietud es una de las razones por las que Shannon desarrolló la
Teoría de la Información, que permite \textbf{cuantificar la cantidad de información} de un
mensaje. De allí dos conceptos:

\begin{defbox}{Redundancia y distancia de unicidad}
\begin{itemize}[leftmargin=1.4em,itemsep=2pt]
  \item \textbf{Redundancia}: diferencia entre el \emph{tamaño} de un mensaje y la
        \emph{información} que realmente lleva. Los lenguajes naturales tienen mucha
        redundancia (por eso entendemos una oración a la que le sacaron las vocales). Es un
        \emph{mecanismo de control de errores}: si se pierde una letra, el sentido se
        mantiene.
  \item \textbf{Distancia de unicidad}: longitud mínima que debe tener un mensaje para que las
        chances de que exista \emph{más de un} descifrado con sentido sean despreciables.
\end{itemize}
\end{defbox}

\begin{notebox}
\textbf{Datos prácticos.}
\begin{itemize}[leftmargin=1.3em,itemsep=1pt]
  \item \textbf{Compresión $=$ eliminar redundancia}: comprimir es recodificar en un alfabeto
        donde el tamaño del mensaje iguala al de la información. El \emph{algoritmo de Huffman}
        es una aplicación directa de la teoría de Shannon.
  \item \textbf{Distancia de unicidad de los lenguajes naturales} (español, inglés): del orden
        de \textbf{4 caracteres}. Cualquier mensaje de más de 4 letras tiene chance
        despreciable de tener dos descifrados con sentido. Por eso, en la práctica, la fuerza
        bruta sobre ROT-X funciona: el primer descifrado con sentido suele ser el correcto.
\end{itemize}
\end{notebox}

\begin{warnbox}{Los límites de la fuerza bruta}
La fuerza bruta \emph{siempre} está disponible para el atacante, pero \textbf{tiene un
prerrequisito}: poder \emph{discriminar un descifrado válido de uno inválido}. \\[3pt]
\textbf{Contraejemplo (caso límite).} Supongamos que se cifró con ROT-X un mensaje plano $p$
de \textbf{una sola letra}, elegida al azar de una distribución \emph{uniforme}, y el texto
cifrado resultante es $c=\texttt{a}$. ¿Cuál era $p$? \textbf{Imposible de saber}: cualquier
letra es un descifrado igualmente válido, así que aun sobre un cifrado ``trivialmente
inseguro'' la fuerza bruta no decide. \\[3pt]
\emph{Hipótesis de fuerza bruta}: se puede discriminar un descifrado válido de uno inválido.
Esta idea se retoma al final de la clase (es la base de uno de los desarrollos seminales de la
criptografía moderna).
\end{warnbox}

%=====================================================================
\section{Cifrado de sustitución (monoalfabético)}
%=====================================================================

\begin{defbox}{Cifrado de sustitución}
Para cada símbolo del alfabeto se establece una \textbf{transformación única} (se reemplaza
cada símbolo por otro). La asignación debe ser \emph{biyectiva}: \textbf{sin repeticiones}
(dos símbolos distintos no pueden ir al mismo, o no se podría invertir). Los símbolos de
llegada \emph{no} tienen por qué ser del mismo alfabeto: pueden ser otros símbolos.
\end{defbox}

\begin{practbox}{Ejemplo trabajado}
Clave (permutación del alfabeto):
\begin{center}\ttfamily
\begin{tabular}{ll}
k = & dublcmfthijnzpxqeaosvkrwgy\\
    & abcdefghijklmnopqrstuvwxyz
\end{tabular}
\end{center}
La fila de abajo es el alfabeto plano; la de arriba, su imagen. Entonces:
\[
  \texttt{esto es una prueba}\;\xrightarrow{\;\text{sustitución}\;}\;\texttt{cosx co vpd qavcud}
\]
(cada \texttt{e} se transforma siempre en \texttt{c}, cada \texttt{s} en \texttt{o}, etc.)
\end{practbox}

\subsection{Tamaño del espacio de claves}
\begin{practbox}{¿Cuántas claves tiene? — cálculo}
La clave es la \emph{transformación} completa. Como no puede haber repeticiones, quedarse con
la fila imagen equivale a una \textbf{permutación} del alfabeto. Para el español
(\textbf{27 símbolos}):
\[
  |\Kset| \;=\; 27! \;\approx\; 1.09\times10^{28}
  \qquad(\text{un número de} \approx 28 \text{ dígitos}).
\]
Es un número ``bestialmente alto'' (del orden de cuatrillones, en escala larga). Por lo tanto,
a diferencia de la rotación, \textbf{este cifrado NO tiene pocas claves}: la fuerza bruta es
inviable.
\end{practbox}

\subsection{Criptoanálisis: análisis de frecuencias}
A pesar del enorme espacio de claves, el cifrado de sustitución es inseguro.

\begin{warnbox}{Debilidad: las propiedades estadísticas del lenguaje no se alteran}
La transformación \textbf{preserva las características estadísticas del lenguaje}. Como las
vocales aparecen mucho más que las consonantes, si en el texto cifrado un símbolo aparece
mucho, hay altas chances de que sea la imagen de una vocal. La misma letra del plano produce
\emph{siempre} la misma letra cifrada, lo que delata patrones (p.\ ej.\ en ``Eduardo'' la 2.\textsuperscript{a}
y la anteúltima letra cifradas serán iguales).
\end{warnbox}

\begin{defbox}{Procedimiento de criptoanálisis}
\begin{enumerate}[leftmargin=1.6em,itemsep=1pt]
  \item Obtener la frecuencia \emph{estimada} de cada símbolo en el lenguaje del mensaje
        (\emph{tabla de frecuencias} del idioma).
  \item Calcular la frecuencia de cada símbolo en el \emph{texto cifrado}.
  \item Asumir que los símbolos de mayor probabilidad se corresponden.
  \item Formar grupos de dos y tres letras comunes (\texttt{el}, \texttt{la}, \texttt{de},
        \texttt{las}, \texttt{los}, \dots) y propagar/podar por prueba y error.
\end{enumerate}
\end{defbox}

% ---- Diagrama: gráfico de barras de frecuencias (slide Frecuencias de letras) ----
\begin{figure}[h]
\centering
\begin{tikzpicture}[font=\scriptsize]
  \def\sx{0.42} % ancho de banda
  \def\sy{0.28} % escala vertical (unidades por %)
  % ejes
  \draw[->] (-0.4,0) -- (27.5*\sx+0.4,0) node[right]{letra};
  \draw[->] (-0.4,0) -- (0,14*\sy+0.4) node[above]{\% frecuencia};
  \foreach \v in {0,5,10}{
     \draw[black!30] (-0.1,\v*\sy) -- (27*\sx,\v*\sy);
     \node[anchor=east] at (-0.15,\v*\sy) {\v};
  }
  % datos: letra/frecuencia (tabla del español de la slide)
  \foreach \i/\L/\f in {
     0/E/13.11, 1/A/10.60, 2/S/8.47, 3/O/8.23, 4/I/7.16, 5/N/7.14,
     6/R/6.95, 7/D/5.87, 8/T/5.40, 9/C/4.85, 10/L/4.42, 11/U/4.34,
     12/M/3.11, 13/P/2.71, 14/G/1.40, 15/B/1.16, 16/F/1.13, 17/V/0.82,
     18/Y/0.79, 19/Q/0.74, 20/H/0.60, 21/Z/0.26, 22/J/0.25, 23/X/0.15,
     24/W/0.12, 25/K/0.11, 26/{Ñ}/0.10}{
     \fill[histBar] ({\i*\sx+0.07},0) rectangle ({\i*\sx+\sx-0.07},{\f*\sy});
     \node[anchor=north,font=\tiny] at ({\i*\sx+\sx/2},-0.05) {\L};
  }
\end{tikzpicture}
\caption{Tabla de frecuencias del idioma español (en \%). Frecuencia alta: E, A, S, O, I, N,
R, D, T; media: C, L, U, M, P, G, B, F; baja: V, Y, Q, H, Z, J, X, W, K, Ñ.
\textit{(Recreación, como gráfico de barras, de la diapositiva ``Frecuencias de letras''.)}}
\end{figure}

\begin{notebox}
\textbf{Tabla de frecuencias del español} (valores de la slide, en \%):
\begin{center}\ttfamily\footnotesize
\begin{tabular}{llllll}
E 13,11 & C 4,85 & Y 0,79 \qquad & A 10,60 & L 4,42 & Q 0,74\\
S 8,47  & U 4,34 & H 0,60 \qquad & O 8,23  & M 3,11 & Z 0,26\\
I 7,16  & P 2,71 & J 0,25 \qquad & N 7,14  & G 1,40 & X 0,15\\
R 6,95  & B 1,16 & W 0,12 \qquad & D 5,87  & F 1,13 & K 0,11\\
T 5,40  & V 0,82 & Ñ 0,10 \qquad &         &        &      \\
\end{tabular}
\end{center}
La tabla es una característica \emph{intrínseca} del lenguaje: si se calcula sobre 3 o 4 libros
grandes, el histograma casi no varía.
\end{notebox}

\begin{notebox}
\textbf{Anécdotas del profesor.}
\begin{itemize}[leftmargin=1.3em,itemsep=1pt]
  \item Las revistas de \emph{acertijos / criptogramas} para chicos de primaria son
        exactamente esto: un chico de 8--9 años puede resolver un cifrado de sustitución sin
        teoría, usando su noción estadística internalizada del lenguaje.
  \item Pistas idiomáticas: en español la \texttt{q} casi siempre va seguida de \texttt{u}; tres
        consonantes seguidas no existen; los artículos (\texttt{el}, \texttt{la}, \texttt{de},
        \texttt{los}) aparecen mucho.
  \item En un curso de criptoanálisis, el profesor descifró un mensaje en \textbf{alemán} sin
        que nadie supiera alemán, apoyándose solo en los \emph{indicadores estadísticos
        formales} del idioma. Reveló que el ataque no requiere conocer el lenguaje.
  \item \textbf{Aplica al mundo digital}: cifrar un \texttt{.mp4} o un \texttt{.docx} con
        sustitución es igualmente atacable, porque el formato tiene mucha \emph{estructura}
        (que pasa a ser parte de la estadística del ``lenguaje''). A veces es más fácil, porque
        los \emph{headers} deben tener valores concretos.
\end{itemize}
\end{notebox}

\subsection{Los símbolos pueden ser diferentes: el criptograma de Trinity}
\begin{notebox}
\textbf{Criptograma del cementerio de Trinity (NY, 1794).} Ejemplo de sustitución donde cada
letra se reemplaza por un \emph{símbolo} (no por otra letra). La clave se arma con
\textbf{tableros tipo ta-te-ti}: como hay más de un tablero, se los distingue con cero, uno o
dos puntos, y cada letra se reemplaza por la representación de su cuadrante. Fue descifrado
recién en \textbf{1986}: no por dificultad del sistema, sino porque el \emph{mensaje era muy
corto} (poca información para el ataque estadístico). \emph{(Queda de tarea descifrarlo.)}
\end{notebox}

\subsection{Ejercicio resuelto en clase: ``el escarabajo de oro''}
\begin{practbox}{Ejercicio de descifrado (hecho en vivo en la clase)}
Se planteó descifrar un mensaje en \textbf{castellano} cifrado por sustitución con
\emph{símbolos / números}. Ayudas dadas:
\begin{itemize}[leftmargin=1.3em,itemsep=1pt]
  \item El mensaje original está en castellano.
  \item La separación en grupos de 5 símbolos no es parte del problema; solo ayuda a contar.
  \item \textbf{Gancho} (\emph{crib}): la secuencia \texttt{LACABEZA} aparece en el mensaje
        plano.
\end{itemize}
El mensaje cifrado (con palos de cartas, llaves y dígitos) fue, por ejemplo:
\begin{center}\ttfamily\small
$\diamondsuit$>]$\clubsuit$0\quad $\heartsuit$34$\diamondsuit$3\quad 44]>\}\quad
>\{:$\odot$$\spadesuit$\quad \{4$\heartsuit$3>\quad 34]02\quad 014:0\quad \dots
\end{center}
\textbf{No se dio la solución}: el mensaje pertenece a \emph{El escarabajo de oro}, de
\textbf{Edgar Allan Poe} (también criptógrafo amateur), cuyo libro explica cómo decodificarlo;
quien lo termine puede verificar por ahí.
\end{practbox}

\begin{warnbox}{Lección: ``débil'' no es lo mismo que ``romperlo''}
Una cosa es decir, desde lo abstracto, que un criptosistema es débil/fácil de romper; otra es
\emph{romperlo}. Cuanto más complejo el criptosistema, más grande el espacio de búsqueda.
\textbf{Contracara histórica}: hubo criptosistemas considerados ``buenos'' \emph{solo porque
nadie los había roto} —no por ser fundacionalmente seguros, sino porque nadie les había
dedicado el tiempo necesario—. Esta es la idea de ``\emph{handshake / crib}'' (gancho) que se
retoma con Enigma.
\end{warnbox}

%=====================================================================
\section{Cifrado de Vigenère (sustitución polialfabética)}
%=====================================================================

Al entrar el Renacimiento surgen los \textbf{criptosistemas de sustitución polialfabética}:
también sustituyen símbolo por símbolo según una regla, pero con la característica de que
\textbf{un mismo símbolo puede terminar en símbolos distintos}. El más emblemático es el de
\textbf{Vigenère}.

\begin{defbox}{Cifrado de Vigenère}
Es una \textbf{composición de cifrados por rotación}. Se define una clave de longitud $n$,
donde cada elemento indica \emph{cuántas veces rotar} la letra en esa posición. El mensaje se
divide en \textbf{bloques} del tamaño de la clave; al terminar la clave, se vuelve a empezar
(la clave se repite cíclicamente).
\end{defbox}

\begin{practbox}{Ejemplo trabajado: $K=\texttt{ECFD}=4253$ (long.\ 4)}
La clave puede darse como letras o como números (con la correspondencia
$\texttt{A}=0,\texttt{B}=1,\texttt{C}=2,\dots$; aquí se usa la variante numérica de la slide).
Se rota cada letra según el dígito de la clave alineado debajo:
\begin{center}\ttfamily\small
\begin{tabular}{ll}
mensaje:  & esto es una prueba\\
clave:    & 4253 42 534 253 42\\
cifrado:  & iuyr dg xcy ofc ttzhfc
\end{tabular}
\end{center}
\textbf{Lo clave del ejemplo:} la letra \texttt{e}, según en qué posición del bloque caiga, se
transforma en \emph{letras distintas}; y una misma letra del cifrado (p.\ ej.\ dos
\texttt{c} seguidas, o un espacio que aparece dos veces) puede provenir de letras \emph{distintas}
del plano. Esto \textbf{rompe todas las propiedades estadísticas}: la letra más frecuente del
plano deja de ser el símbolo más frecuente del cifrado.
\end{practbox}

\begin{notebox}
\textbf{Vigenère en la historia.} Creado en \textbf{1553}. Durante casi \textbf{300 años} se
lo consideró seguro. En \textbf{1863}, \textbf{Friedrich Kasiski} (un soldado aficionado a la
criptografía, destacado en la fría Rusia y con mucho tiempo libre) publica un método para
resolverlo.
\end{notebox}

\subsection{Criptoanálisis: el test de Kasiski}
El ataque tiene \textbf{dos fases}: (1) determinar la longitud de la clave (tamaño de bloque);
(2) analizar la clave de cada bloque por separado. Una consecuencia importante: para los
polialfabéticos \textbf{ya no hay solución genérica}; los ataques pasan a ser
\emph{ad hoc}, específicos de cada algoritmo (explotan el conocimiento del criptosistema).

\subsubsection{Fase 1 — Longitud de la clave (n-gramas repetidos)}
\begin{defbox}{Observación de Kasiski}
Si una misma secuencia del texto plano cae en la \emph{misma posición de bloque}, se cifra
igual. Por lo tanto, \textbf{secuencias repetidas en el texto cifrado} provienen, con alta
probabilidad, de los mismos caracteres del plano alineados al mismo punto del bloque. La
\textbf{distancia} entre dos apariciones debe ser un \textbf{múltiplo del tamaño de bloque}.
\[
  \text{longitud de clave} \;=\; \mcd(\text{distancias entre repeticiones})
\]
\end{defbox}

\begin{practbox}{Ejemplo trabajado del test de Kasiski (slide)}
Sobre un texto cifrado largo se hallan estas secuencias repetidas y sus distancias:
\begin{center}
\renewcommand{\arraystretch}{1.15}
\begin{tabular}{lll}
\toprule
secuencia & apariciones & distancias\\
\midrule
\texttt{GGMP} & 3 veces & 256 y 104 caracteres\\
\texttt{YEDS} & 2 veces & 72 caracteres\\
\texttt{HASE} & 2 veces & 156 caracteres\\
\texttt{VSUE} & 2 veces & 32 caracteres\\
\bottomrule
\end{tabular}
\end{center}
Cálculo del máximo común divisor:
\[
  \mcd(32,\;72,\;104,\;156,\;256) \;=\; 4 .
\]
Luego, la \textbf{longitud de clave es 4}. (Es una prueba estadística: si falla, suele dar
$\mcd=1$ por una coincidencia espuria; se descarta esa muestra y se recalcula.)
\end{practbox}

\subsubsection{Fase 2 — Método de coincidencia mutua}
Conocida la longitud $n$, el mensaje se separa en $n$ \emph{sub-mensajes}: el formado por las
primeras letras de cada bloque, el de las segundas, etc. Cada sub-mensaje fue cifrado con
\textbf{una única rotación} (es, en sí, un cifrado por rotación). Como un histograma de letras
es una propiedad robusta del lenguaje, cada sub-mensaje conserva (rotado) la estadística del
idioma.

\begin{defbox}{Método de coincidencia mutua}
\begin{enumerate}[leftmargin=1.6em,itemsep=1pt]
  \item Determinar estadísticamente las letras más probables en cada grupo.
  \item Estimar las rotaciones que llevan a cada grupo a \emph{equipararse} con las
        estadísticas del lenguaje.
  \item Probar las rotaciones, podando caminos al encontrar combinaciones sin sentido.
\end{enumerate}
\end{defbox}

\begin{notebox}
\textbf{La idea fina de Kasiski.} Si se juntan dos sub-grupos con rotaciones distintas, el
histograma combinado ``se rompe''. Pero \emph{existe una rotación relativa} entre ambos para
la cual, al combinarlos, el histograma vuelve a parecerse al del idioma: esa rotación es la
\textbf{diferencia} entre las dos claves parciales. Encontrando todas las diferencias
relativas, se transforma el problema en un \emph{único} cifrado por rotación, y entonces se
prueban las $\sim$26 posibilidades para recuperar el mensaje.
\end{notebox}

%=====================================================================
\section{Sustitución polialfabética avanzada: la máquina Enigma}
%=====================================================================

El estado del arte hacia la Segunda Guerra Mundial eran polialfabéticos mucho más
sofisticados que Vigenère. Todos los ejércitos (Alemania, EE.UU., Japón, Inglaterra) tenían
alguno. El más estudiado es \textbf{Enigma}, usado por los alemanes.

\begin{defbox}{Enigma — criptosistema de rotores}
Es electromecánica, similar a una máquina de escribir, salvo que al pulsar una tecla
\emph{imprime otra letra}. Componentes:
\begin{itemize}[leftmargin=1.4em,itemsep=2pt]
  \item \textbf{Rotores} (3, en versiones posteriores 4): cada rotor es una sustitución
        simple (entra una letra, sale otra). Al \emph{rotar}, la asociación cambia.
  \item \textbf{Avance escalonado}: cada vez que se pulsa una tecla, el primer rotor avanza;
        al completar una vuelta hace avanzar al segundo, y así (como un cuentakilómetros). El
        período antes de repetir es $\approx 26^3$ posiciones: en términos prácticos, para un
        mensaje \textbf{nunca se repite}.
  \item \textbf{Plugboard} (tablero de conexiones): pares de cables intercambiables.
  \item \textbf{Reflector}: la señal ``rebota'' y vuelve. Gracias a esto, con la misma
        configuración inicial, tipear el cifrado devuelve el plano: la \textbf{misma máquina
        cifra y descifra} (idea reutilizada en sistemas modernos).
\end{itemize}
\end{defbox}

% ---- Diagrama: esquema de Enigma (slide Maquina enigma) ----
\begin{figure}[h]
\centering
\begin{tikzpicture}[font=\scriptsize,>=Stealth]
  \tikzset{rotor/.style={draw,fill=itbaBlue!40,minimum width=1.0cm,minimum height=3.0cm}}
  % teclado
  \node[draw,fill=black!8,minimum width=0.9cm,minimum height=2.6cm,align=center] (kb) at (0,0)
       {\rotatebox{90}{Teclado}};
  % rotores
  \node[rotor] (r1) at (1.9,0) {26};
  \node[rotor] (r2) at (3.1,0) {26};
  \node[rotor] (r3) at (4.3,0) {26};
  \node[draw,fill=black!25,minimum width=0.45cm,minimum height=3.0cm] (refl) at (5.1,0) {};
  \node[align=center,font=\tiny] at (5.1,2.0) {Reflector};
  % plugboard
  \node[draw,fill=black!12,minimum width=0.9cm,minimum height=2.6cm,align=center] (pb) at (6.6,0)
       {\rotatebox{90}{Plugboard}};
  % lamparas
  \node[draw,circle,minimum size=0.45cm] (l1) at (8.1,0.7) {T};
  \node[draw,circle,minimum size=0.45cm,fill=yellow!70] (l2) at (8.1,0.0) {R};
  \node[draw,circle,minimum size=0.45cm] (l3) at (8.1,-0.7) {E};
  \node[above=1mm of l1,font=\tiny] {Lámparas};

  \node[below=1mm of r1] {Rotor 1}; \node[below=1mm of r2] {Rotor 2};
  \node[below=1mm of r3] {Rotor 3};
  % recorrido de la señal (ida: hacia el reflector)
  \draw[->,red,line width=0.9pt]
        (0.45,0.9) -- (1.4,0.9) -- (1.4,0.3) -- (2.5,0.3) -- (2.5,-0.4)
        -- (3.7,-0.4) -- (3.7,0.5) -- (4.9,0.5) -- (4.9,-0.6) -- (4.3,-0.6);
  % recorrido de la señal (vuelta: del reflector a la lampara)
  \draw[->,blue,line width=0.9pt]
        (4.3,-0.9) -- (1.9,-0.9) -- (1.9,-1.3) -- (6.1,-1.3)
        -- (6.1,0.0) -- (7.6,0.0) -- (l2.west);
  \node[red,font=\tiny] at (0.2,1.15) {T};
  \node[blue,font=\tiny] at (8.1,-0.95) {R};
\end{tikzpicture}
\caption{Esquema conceptual de Enigma: la pulsación pasa por el plugboard, atraviesa los 3
rotores, se refleja y regresa por otro camino hasta encender una lámpara. Cada pulsación hace
avanzar un rotor, por lo que la misma tecla produce salidas distintas (``3 permutaciones por
pasada''). \textit{(Recreación de la diapositiva ``Máquina Enigma''.)}}
\end{figure}

\begin{notebox}
\textbf{La clave de Enigma y su criptoanálisis.}
\begin{itemize}[leftmargin=1.3em,itemsep=1pt]
  \item La \textbf{clave} era la \emph{disposición inicial de los rotores} (p.\ ej.\
        $15,8,14$). El plugboard quedaba fijado por tipo de unidad (marina, infantería). Cada
        oficial tenía un \emph{cuaderno} con una configuración distinta \textbf{por día}: si
        en un día no se barrían todas las posibilidades, el trabajo se perdía.
  \item Los aliados, liderados por \textbf{Alan Turing} (entre otros) en \emph{Bletchley
        Park}, \textbf{capturaron físicamente} una máquina, hicieron \emph{ingeniería
        inversa} y construyeron una computadora (la \emph{Bombe}) para fuerza bruta. Aun así,
        $26^3$ claves por día desbordaban el cómputo.
  \item Lo lograron usando \textbf{cribs} (ganchos): muchos mensajes empezaban con el mismo
        \emph{saludo}, lo que reducía drásticamente el espacio. \emph{No rompieron Enigma de
        forma general}, solo ese \textbf{uso} de Enigma. Los submarinos con máquinas de
        \textbf{4 rotores} no pudieron descifrarse durante la guerra.
  \item \textbf{Legado}: para resolver esto prácticamente se inventó la computadora; de aquí
        sale el modelo de procesamiento que luego von Neumann toma para la arquitectura
        (memoria + ejecución) que usamos hoy.
\end{itemize}
\end{notebox}

%=====================================================================
\section{¿Hay criptosistemas seguros? El nacimiento de la criptografía moderna}
%=====================================================================

La historia de los clásicos es un juego del gato y el ratón: se inventan criptosistemas, se
rompen, se inventan otros. La Segunda Guerra Mundial genera una \textbf{crisis}: al juntar
mucha gente brillante, se rompieron cosas que se suponían muy seguras. La realización clave
fue que \textbf{no había un criterio} para dar fiabilidad: ``esto es seguro porque nadie lo
rompió'' no es una definición.

\begin{defbox}{Los tres pilares de la criptografía moderna}
\begin{enumerate}[leftmargin=1.6em,itemsep=2pt]
  \item \textbf{Representaciones formales} de criptosistemas (ya no recetas algorítmicas, sino
        definiciones precisas).
  \item \textbf{Definiciones del modelo de amenaza}: \emph{no hay una única definición de ``es
        seguro''}. (Tal vez Enigma era seguro, pero el \emph{uso} que se le daba no lo era.)
  \item \textbf{Pruebas formales de seguridad}, que permiten salir de la definición vaga y
        \emph{comparar} la seguridad de dos construcciones.
\end{enumerate}
\end{defbox}

\subsection{Definición formal de criptosistema}
\begin{defbox}{Criptosistema (terna de algoritmos)}
\[
\begin{aligned}
  \Gen&: () \to \Kset && \text{Generador de clave}\\
  \Enc&: \Kset \times \Pset \to \Cset && \text{Cifrado}\\
  \Dec&: \Kset \times \Cset \to \Pset && \text{Descifrado}
\end{aligned}
\]
\textbf{Propiedad básica (corrección).} Para todo $m$ y $k$ válidos:
\[
  d_k\bigl(e_k(m)\bigr) = m .
\]
\textbf{Conjuntos involucrados:} $\Kset$ (espacio de claves: todas las claves posibles),
$\Pset$ (espacio plano: todos los mensajes posibles), $\Cset$ (espacio cifrado: todos los
mensajes cifrados posibles).
\end{defbox}

\begin{notebox}
\textbf{Detalles.}
\begin{itemize}[leftmargin=1.3em,itemsep=1pt]
  \item $\Gen$ suele ser una función \emph{sin parámetros} que \textbf{determina el espacio de
        claves} (dice cuántas y cuáles claves hay: 26, un millón, un cuatrillón, \dots).
  \item $\Enc$ y $\Dec$ deben actuar como \textbf{inversas} para una misma clave: si cifro con
        $k$ y descifro con $k$, recupero el mensaje.
  \item Esta propiedad es \emph{``sanitaria''} (básica): garantiza que el sistema sirva, pero
        \textbf{no habla de seguridad}. De hecho, la función \emph{identidad} (que ignora la
        clave y devuelve el mensaje) la cumple. No se incluye una propiedad de seguridad
        \emph{porque no hay una única noción de seguridad}.
\end{itemize}
\end{notebox}

%=====================================================================
\section{Seguridad: el secreto perfecto de Shannon}
%=====================================================================

Al formalizar la seguridad, \textbf{Shannon} define el criterio de \textbf{secreto perfecto}:
una formulación estadística que constituye el \emph{ideal} de los criptosistemas.

\begin{defbox}{Secreto perfecto (Shannon)}
Dado un criptosistema $(\Gen,e,d)$, posee \textbf{secreto perfecto} si para toda distribución
de probabilidades en $\Mset$, para cada mensaje $m$ y cada cifrado $c$ con $\Pr[C=c]>0$:
\[
  \Pr[\,M=m \mid C=c\,] \;=\; \Pr[\,M=m\,] .
\]
Es decir: \textbf{conocer el texto cifrado no modifica} la probabilidad \emph{a priori} que se
tenía sobre el mensaje. No se aprende \emph{nada} del mensaje a partir del cifrado.
\end{defbox}

\begin{notebox}
\textbf{Interpretación (con ejemplos del profesor).}
\begin{itemize}[leftmargin=1.3em,itemsep=1pt]
  \item \textbf{Ejemplo de una letra}: si la probabilidad \emph{a priori} de que el mensaje
        sea \texttt{a} es $0.106$ (su frecuencia), y tras conocer el cifrado $c=\texttt{J}$
        sigue siendo $0.106$ (y lo mismo para todas las letras), el sistema tiene secreto
        perfecto.
  \item \textbf{Ejemplo binario}: si la respuesta es ``sí/no'' y se sabe a priori que hay
        80\,\% de chances de ``sí'', con secreto perfecto conocer el cifrado mantiene ese
        80\,\%: no pasa a 79\,\% ni a 81\,\%. \emph{Es mucho más fuerte que ``no poder
        decodificar''}: es no extraer \emph{ninguna} información.
\end{itemize}
\end{notebox}

\begin{warnbox}{La paradoja del secreto perfecto}
Notar que $c=e_k(m)$, así que las variables aleatorias discretas $C$ y $M$ están
\emph{relacionadas por diseño} (la función de cifrado). Sin embargo, la propiedad de secreto
perfecto, interpretada probabilísticamente, dice que $M$ y $C$ se comportan como variables
\textbf{estadísticamente independientes} (conocer una no modifica la distribución de la otra).
Esa tensión es la ``paradoja''. \\[3pt]
\emph{Consecuencia bestial:} un criptosistema con secreto perfecto es \textbf{imposible de
romper, incluso por fuerza bruta} (conecta con el límite de la fuerza bruta visto en ROT-X).
\end{warnbox}

\begin{notebox}
\textbf{Pregunta abierta que deja la clase.} ¿El secreto perfecto es una \emph{asíntota
teórica} inalcanzable (que solo sirve para medir cuán cerca estamos), o existe una
construcción que efectivamente lo alcanza? Se retoma la clase siguiente, junto con un repaso
de probabilidades y los matices que llevan a distintas definiciones de seguridad.
\end{notebox}

%=====================================================================
\section*{Lectura recomendada}
\addcontentsline{toc}{section}{Lectura recomendada}
%=====================================================================
\begin{center}
\begin{tcolorbox}[enhanced,colback=itbaBlue!20,colframe=itbaCyanDark,boxrule=1pt,
  arc=3pt,width=0.85\textwidth,halign=center]
\large\textbf{Capítulo 1}\\[4pt]
\emph{Introduction to Modern Cryptography}\\[2pt]
J.~Katz \& Y.~Lindell
\end{tcolorbox}
\end{center}
\noindent Recoge todo lo visto y contiene más ejemplos de criptosistemas clásicos.

%---------------------------------------------------------------------
\vfill
\begin{center}\small\itshape\color{black!60}
Apunte integral de la Clase 1 — Introducción a la Criptografía · Criptografía y Seguridad, ITBA.\\
Síntesis de la teoría (transcripción + diapositivas). Esta clase no tuvo clase práctica asociada.
\end{center}

\end{document}
